% ============================================================================
% CAPÍTULO 3 — CSS\index{CSS} moderno: layout, responsividade\index{responsividade} e estilo
% ============================================================================
\chapter{CSS Moderno: Layout, Responsividade e Estilo}
\label{cap:css}

\begin{caixaconceito}
\textbf{CSS} (Cascading Style Sheets) é a linguagem de \emph{apresentação} da
web. Se o HTML é o esqueleto, o CSS é a pele, a roupa e o cenário: cores,
tipografia, espaçamento, posicionamento e comportamento visual.
\end{caixaconceito}

Ao final deste capítulo, você será capaz de:
\begin{objetivos}
  \item Aplicar o modelo de cascata\index{cascata}, herança e especificidade\index{especificidade} sem sustos;
  \item Dominar o \emph{box model} e o \texttt{box-sizing};
  \item Construir layouts com \textbf{Flexbox} e \textbf{Grid};
  \item Criar páginas responsivas com \emph{mobile-first} e media queries\index{media queries};
  \item Organizar CSS com variáveis, BEM e boas práticas;
  \item Entregar o projeto do capítulo: uma landing page responsiva.
\end{objetivos}

\section{Como o CSS funciona: cascata, herança e especificidade}

Três forças decidem qual regra vence quando duas regras conflitam:

\begin{enumerate}[leftmargin=*, itemsep=0.4em]
  \item \textbf{Especificidade}: quanto mais específico o seletor, maior o peso.
  \texttt{\#id} $>$ \texttt{.classe} $>$ \texttt{tag}. Na prática, 99\% do tempo
  você usará classes;
  \item \textbf{Cascata}: com a mesma especificidade, vence a regra declarada
  por último no arquivo;
  \item \textbf{Herança}: propriedades de texto (cor, fonte) são herdadas pelos
  filhos; propriedades de layout (margin, padding) não são.
\end{enumerate}

\begin{lstlisting}[style=livro, language=CSS,
  caption={Cascata e especificidade na prática.}, label={list:css-cascata}]
/* 1. Regra por tag — peso baixo */
button { background-color: #cccccc; }

/* 2. Regra por classe — peso maior */
.botao-primario { background-color: #1f4e79; }

/* 3. Regra com !important — só em último caso! */
button { background-color: red !important; }
\end{lstlisting}

\begin{caixaarmadilha}
\texttt{!important} quebra a cascata e vira uma bomba-relógio de manutenção.
Use-o apenas para sobrepor estilos de bibliotecas de terceiros — e mesmo assim,
prefira especificidade bem escrita.
\end{caixaarmadilha}

\section{O box model: tudo é uma caixa}

Todo elemento HTML é uma caixa com quatro camadas (de dentro para fora):
\textbf{conteúdo} $\rightarrow$ \textbf{padding} (respiro interno)
$\rightarrow$ \textbf{border} $\rightarrow$ \textbf{margin} (espaço externo).

\begin{lstlisting}[style=livro, language=CSS,
  caption={Box model com \texttt{box-sizing: border-box}.}, label={list:css-box}]
* {
  box-sizing: border-box; /* largura inclui padding e borda */
}

.cartao {
  width: 300px;
  padding: 16px;
  border: 1px solid #ddd;
  margin: 12px auto;
}
\end{lstlisting}

\begin{caixaconceito}
Com \texttt{box-sizing: border-box}, \texttt{width: 300px} é o tamanho
\emph{final} da caixa (conteúdo + padding + borda). Sem ele, o CSS soma tudo e
as caixas ``estouram'' — a causa nº 1 de layouts quebrados em iniciantes.
O reset universal (\texttt{* \{ box-sizing: border-box; \}}) é padrão de mercado.
\end{caixaconceito}

\section{Flexbox: alinhamento em uma dimensão}

Flexbox\index{Flexbox} organiza itens em \emph{uma linha ou uma coluna} e resolve 90\% dos
problemas de alinhamento do dia a dia.

\begin{lstlisting}[style=livro, language=CSS,
  caption={Flexbox: centralizar e distribuir.}, label={list:css-flex}]
.navbar {
  display: flex;
  justify-content: space-between; /* eixo principal */
  align-items: center;            /* eixo transversal */
  gap: 1rem;
}

.menu {
  display: flex;
  gap: 1.5rem;
  list-style: none;
}

.centro {
  display: flex;
  justify-content: center;
  align-items: center;
  min-height: 100vh;
}
\end{lstlisting}

\begin{caixadica}
Os três comandos que você mais vai digitar na vida: \texttt{display: flex},
\texttt{justify-content} (eixo principal) e \texttt{align-items} (eixo
transversal). Memorize o \textbf{jogo Flexbox Froggy}
(\url{https://flexboxfroggy.com}) para internalizar de vez.
\end{caixadica}

\section{Grid: layout em duas dimensões}

Grid\index{Grid} organiza itens em \emph{linhas e colunas simultaneamente} — perfeito para
a estrutura geral da página (header, sidebar, main, footer) e para grades de
cards.

\begin{lstlisting}[style=livro, language=CSS,
  caption={Grid responsivo sem media query (auto-fit + minmax).},
  label={list:css-grid}]
.grade-produtos {
  display: grid;
  grid-template-columns: repeat(auto-fit, minmax(220px, 1fr));
  gap: 1.5rem;
}

.layout-pagina {
  display: grid;
  grid-template-columns: 240px 1fr;
  grid-template-areas:
    "header header"
    "sidebar main"
    "footer footer";
}
\end{lstlisting}

\begin{caixaconceito}
\texttt{repeat(auto-fit, minmax(220px, 1fr))} cria colunas de no mínimo 220px
que crescem igualmente e \emph{quebram} automaticamente conforme a tela
encolhe — responsividade ``de graça'', sem uma única media query.
\end{caixaconceito}

\section{Cores, tipografia e espaçamento}

\begin{lstlisting}[style=livro, language=CSS,
  caption={Sistema de design com variáveis CSS.}, label={list:css-variaveis}]
:root {
  --cor-primaria: #1f4e79;
  --cor-secundaria: #2e86ab;
  --cor-texto: #22313f;
  --espaco-pequeno: 0.5rem;
  --espaco-medio: 1rem;
  --raio: 8px;
  --fonte-corpo: "Inter", system-ui, sans-serif;
}

body {
  font-family: var(--fonte-corpo);
  color: var(--cor-texto);
  line-height: 1.6;
}

.cabecalho {
  background-color: var(--cor-primaria);
  padding: var(--espaco-medio);
  border-radius: var(--raio);
}
\end{lstlisting}

\begin{caixadica}
Use \textbf{unidades relativas}: \texttt{rem} (baseado no tamanho da raiz) para
fontes, \texttt{em} para propriedades que devem escalar com o elemento,
\texttt{\%} e \texttt{fr} para layouts. Evite \texttt{px} em fontes e evite
\texttt{pt} em telas. Isso garante que o texto respeite o zoom do usuário.
\end{caixadica}

\section{Responsividade: mobile-first}

A regra de ouro: \textbf{escreva o CSS para o celular primeiro} e vá
adicionando melhorias para telas maiores com media queries
(\texttt{min-width}). Isso força decisões simples desde o início e reduz
drasticamente o CSS total.

\begin{lstlisting}[style=livro, language=CSS,
  caption={Estratégia mobile-first com media query.}, label={list:css-media}]
/* Base: celular (uma coluna, empilhado) */
.menu {
  display: flex;
  flex-direction: column;
  gap: 0.5rem;
}

/* A partir de 768px: tabelas e desktops */
@media (min-width: 768px) {
  .menu {
    flex-direction: row;
    justify-content: flex-end;
  }
}
\end{lstlisting}

\begin{caixaarmadilha}
Não fixe larguras (\texttt{width: 900px}) em containers: use
\texttt{max-width: 960px; margin-inline: auto}. Assim a página se adapta a
qualquer tela sem barras de rolagem horizontais — o sintoma clássico de layout
quebrado.
\end{caixaarmadilha}

\section{Pseudo-classes, transições e o toque final}

\begin{lstlisting}[style=livro, language=CSS,
  caption={Interações: hover, foco e transição.}, label={list:css-interacao}]
.botao {
  background-color: var(--cor-primaria);
  color: #ffffff;
  padding: 0.75rem 1.5rem;
  border-radius: 6px;
  transition: background-color 200ms ease, transform 200ms ease;
}

.botao:hover {
  background-color: var(--cor-secundaria);
}

.botao:focus-visible {
  outline: 3px solid var(--cor-secundaria);
  outline-offset: 2px;
}

.card:nth-child(even) {
  background-color: #f7f9fb;
}
\end{lstlisting}

\begin{caixaconceito}
\texttt{:hover} é para mouse, \texttt{:focus} é para teclado. Use sempre
\texttt{:focus-visible} para anéis de foco — ele só aparece quando a navegação
é por teclado, sem poluir o visual de quem usa mouse. Nunca remova o outline
sem fornecer um substituto acessível.
\end{caixaconceito}

\section{Arquitetura CSS: organização que escala}

\begin{itemize}[leftmargin=*, itemsep=0.4em]
  \item \textbf{Variáveis} (\texttt{--nome}) no \texttt{:root} centralizam o
  design system;
  \item \textbf{BEM} (Block\_\_Element--Modifier) dá nomes previsíveis:
  \texttt{.botao}, \texttt{.botao\_\_icone}, \texttt{.botao--primario};
  \item \textbf{Arquivos organizados}: \texttt{reset.css}, \texttt{variaveis.css},
  \texttt{componentes/}, \texttt{layout/} — ou um único arquivo em projetos pequenos;
  \item \textbf{Metodologias}: você estudará Tailwind CSS (utility-first) a
  fundo no capítulo~\ref{cap:estilizacao} — a abordagem dominante do mercado em 2026.
\end{itemize}

% ============================================================================
\section{Projeto do capítulo: landing page responsiva}
\label{sec:projeto3}

\begin{caixaprojeto}
\textbf{Objetivo:} criar a landing page de um produto fictício — um app de
finanças pessoais chamado \emph{OrçaFácil} — com hero, seção de recursos em
grid, planos e rodapé.

\textbf{Critérios de aceite:}
\begin{itemize}[leftmargin=*, itemsep=0.2em]
  \item \emph{Mobile-first}: perfeita no celular e no desktop (teste em ambos);
  \item Hero com título, subtítulo, CTA e imagem;
  \item Seção de recursos em Grid (\texttt{auto-fit/minmax});
  \item Seção de planos em Flexbox com destaque no plano principal;
  \item Variáveis CSS para o design system;
  \item Sem barra de rolagem horizontal em nenhuma largura.
\end{itemize}
\end{caixaprojeto}

O HTML e o CSS do projeto estão em \texttt{codigo/cap03-landing/} —
\texttt{index.html} e \texttt{style.css}. O trecho essencial do layout
(Listagem~\ref{list:landing-css}):

\lstinputlisting[style=livro, language=CSS, firstline=1, lastline=40,
  caption={Trecho do layout da landing page (Grid + Flexbox).}, label={list:landing-css}]
  {\codigopath/cap03-landing/style.css}

\begin{caixadica}
Teste a responsividade no DevTools: clique no ícone de dispositivo (Ctrl+Shift+M)
e arraste a largura de 320px até 1440px. Todo desenvolvedor full stack faz esse
teste dezenas de vezes por dia.
\end{caixadica}

\section{Exercícios de fixação}

\begin{exercicios}
  \item Explique, com um exemplo próprio, a diferença entre especificidade e cascata.
  \item O que acontece com um elemento de \texttt{width: 300px} sem
  \texttt{box-sizing: border-box} quando recebe \texttt{padding: 30px}? E com ele?
  \item Centralize um botão no centro exato da tela usando Flexbox.
  \item Crie uma grade de 3 colunas que vira 1 coluna no celular, sem media query.
  \item Escreva o CSS de um cartão com sombra, borda arredondada e transição de elevação no hover.
  \item Qual a diferença entre \texttt{em} e \texttt{rem}? Dê um caso de uso para cada.
\end{exercicios}

\section{Desafios}

\begin{desafios}
  \item \textbf{Barra de progresso}: crie uma barra de progresso animada usando
  \texttt{width} em \texttt{\%} e uma transição.
  \item \textbf{Navbar fixa}: navbar fixa no topo com \texttt{position: sticky},
  sombra ao rolar e links com \texttt{:hover} estilizado.
  \item \textbf{Tema claro/escuro}: use \texttt{prefers-color-scheme} para
  alternar variáveis de cor automaticamente.
  \item \textbf{Desafio Froggy}: complete o Flexbox Froggy e o Grid Garden
  (\url{https://cssgridgarden.com}) até o fim.
\end{desafios}

\section{Exercícios de nível sênior}

\begin{seniores}
  \item \textbf{Comparativo de métodos}: implemente a mesma seção de cards de
  3 formas — Flexbox, Grid e (somente conceitualmente) float — e escreva 8
  linhas sobre quando cada uma é a escolha certa.
  \item \textbf{Contraste e acessibilidade}: meça o contraste do seu landing
  page com a ferramenta do Lighthouse e ajuste até passar em AAA no texto grande.
  \item \textbf{Animações com \texttt{@keyframes}:} adicione uma animação de
  entrada suave ao hero, respeitando \texttt{prefers-reduced-motion}.
\end{seniores}

\section{Armadilhas comuns}

\begin{itemize}[leftmargin=*, itemsep=0.4em]
  \item Esquecer o reset de \texttt{box-sizing} (layouts estouram);
  \item Usar \texttt{px} fixos em fontes (quebra o zoom e a acessibilidade);
  \item Margens negativas e \texttt{position: absolute} como ``gambiarra'' de layout;
  \item Escrever CSS \emph{desktop-first} (fica com 3x mais código);
  \item Esquecer o estado de \texttt{:focus-visible} (teclado fica invisível).
\end{itemize}

\section{Resumo do capítulo}

\begin{itemize}[leftmargin=*, itemsep=0.3em]
  \item CSS = apresentação; cascata + especificidade + herança decidem conflitos;
  \item Box model\index{box model} + \texttt{box-sizing: border-box} = layout previsível;
  \item Flexbox resolve uma dimensão; Grid resolve duas;
  \item \texttt{auto-fit/minmax} = responsividade sem media queries;
  \item Mobile-first com \texttt{min-width} é o padrão da indústria;
  \item Variáveis CSS e BEM organizam projetos que crescem.
\end{itemize}

\section{Referências oficiais para aprofundar}

\begin{itemize}[leftmargin=*, itemsep=0.3em]
  \item MDN — CSS: \url{https://developer.mozilla.org/pt-BR/docs/Web/CSS}
  \item MDN — Guia de Flexbox: \url{https://developer.mozilla.org/en-US/docs/Web/CSS/CSS_flexible_box_layout}
  \item MDN — Guia de Grid: \url{https://developer.mozilla.org/en-US/docs/Web/CSS/CSS_grid_layout}
  \item CSS-Tricks — Guia completo de Flexbox: \url{https://css-tricks.com/snippets/css/a-guide-to-flexbox/}
  \item Flexbox Froggy (prática): \url{https://flexboxfroggy.com}
\end{itemize}
